% Part: computability
% Chapter: computability-theory
% Section: non-comp-set

\documentclass[../../../include/open-logic-section]{subfiles}

\begin{document}

\olfileid[pt]{cmp}{thy}{ncp}
\olsection{Existem conjuntos não computáveis}


Vimos acima que todo conjunto computável é computavelmente enumerável.
A recíproca é verdadeira? O resultado a seguir mostra que, em geral,
não é.

\begin{thm}\ollabel{thm:K-0}
Seja $K_0$ o conjunto
$\Setabs{\tuple{e, x}}{\cfind{e}(x) \fdefined}$. Então $K_0$ é
computavelmente enumerável, mas não é computável.
\end{thm}

\begin{proof}
Para ver que $K_0$ é computavelmente enumerável, observe que ele é o
domínio da função~$f$ definida por
\[
f(z) = \umin{y}{(\len{z} = 2 \land T((z)_0, (z)_1, y))}.
\]
Pois, se $\cfind{e}(x)$ está definida, $f(\tuple{e, x})$ encontra uma
sequência de computação que para; se $\cfind{e}(x)$ é indefinida,
$f(\tuple{e, x})$ também é; e, se $z$ nem sequer codifica um par,
então $f(z)$ também é indefinida.

O fato de que $K_0$ não é computável é simplesmente a indecidibilidade
do problema da parada, \olref[hlt]{thm:halting-problem}.
\end{proof}

O conjunto $K_0$ é o conjunto dos pares $\tuple{e,x}$ tais que
$\cfind{e}(x) \fdefined$, isto é, $\tuple{e,x} \in K_0$ se, e somente
se, $\cfind{e}$ está definida (para) na entrada~$x$; por isso, ele
também é chamado de ``conjunto da parada''. O conjunto
$K = \Setabs{e}{\cfind{e}(e) \fdefined}$ é o ``conjunto da
autoparada''. Ele é frequentemente usado como um conjunto indecidível
canônico.

\begin{thm}\ollabel{thm:K}
O conjunto da autoparada
$K = \Setabs{e}{\cfind{e}(e) \fdefined}$ é !!{c.e.}, mas não é
decidível.
\end{thm}

\begin{proof}
  Suponha que $K$ seja decidível, isto é, que sua função característica
  $\Char{K}$ seja computável. Seja
  \[d(e) = \begin{cases}
  1 & \text{se\/ $\Char{K}(e) = 0$}\\
  \fundefined & \text{caso contrário.}
  \end{cases}
  \] 
  Seja $k$ o índice de~$d$, isto é, $d \simeq \cfind{k}$. Então $d(k)
  \simeq \cfind{k}(k)$. Isso contradiz o fato de que $d(k)
  \fdefined$ se, e somente se, $\cfind{k}(k) \fundefined$, que decorre
  da definição de~$d$.

  $K$ é o domínio de $f(x) = \umin{y}{T(x,x,y)}$ e, portanto, é
  !!{c.e.}
\end{proof}

\end{document}

